% Generated draft: requires mathematical and source review.
\documentclass[11pt]{article}
\usepackage[margin=25mm]{geometry}
\usepackage{fontspec}
\setmainfont{Latin Modern Roman}
\usepackage{amsmath,amssymb,mathtools,unicode-math}
\setmathfont{Latin Modern Math}
\usepackage{enumitem,needspace,xurl,hyperref}
\setlength{\parindent}{0pt}
\setlength{\parskip}{5pt}
\setlength{\emergencystretch}{3em}
\newcommand{\N}{\mathbb{N}}
\newcommand{\Z}{\mathbb{Z}}
\newcommand{\Q}{\mathbb{Q}}
\newcommand{\R}{\mathbb{R}}
\newcommand{\C}{\mathbb{C}}
\newcommand{\F}{\mathbb{F}}
\newcommand{\A}{\mathbb{A}}
\newcommand{\PP}{\mathbb{P}}
\newcommand{\E}{\mathbb{E}}
\newcommand{\eps}{\varepsilon}
\newcommand{\dd}{\,\mathrm{d}}
\begin{document}
\section*{UnsolvedMath ID: 30006260}
\textbf{Source problem code:} catalogue entry
\subsection{Definitions and mathematical statement}
The supplied formulation presupposes a base field \(k\), but does not specify that field, its characteristic, or any algebraic-closure or perfectness hypothesis. Relative to the intended base, let \(G\) be a simple adjoint algebraic group and let \(P\subseteq G\) be a nonreduced parabolic subgroup scheme in the sense intended by the source. Here nonreduced means that the structure sheaf has nonzero nilpotent sections. The precise convention for a nonreduced parabolic subgroup scheme, and whether simplicity is meant over \(k\) or geometrically, remain unspecified. Write
\[
X=G/P
\]
for the intended homogeneous quotient, with its left \(G\)-action.

When the automorphism functor is represented by a group scheme \(\operatorname{Aut}_{X/k}\), its value on a \(k\)-scheme \(T\) is the group of \(T\)-automorphisms of \(X\times_k T\). Let \(\operatorname{Aut}_{X/k}^{0}\) denote its identity component, retaining its scheme structure. The supplied statement describes the homomorphism induced by the left action as a natural inclusion
\[
\iota:G\hookrightarrow\operatorname{Aut}_{X/k}^{0}.
\]
The missing setup must justify representability, factorization through the identity component, and the asserted inclusion; faithfulness as a group-scheme action cannot be inferred here merely from the word adjoint.

\textbf{Problem, subject to the unresolved setup.} For the intended class of pairs \((G,P)\), determine necessary and sufficient conditions for \(\iota\) to be an isomorphism of group schemes. In the cases where the asserted inclusion is strict, determine the structure of \(\operatorname{Aut}_{X/k}^{0}\). This second request occurs in the supplied cleaned statement but is not explicit in the supplied original fragment.

The original fragment also asserts that any parabolic subgroup is an intersection of parabolics with maximal reduced part. Its ambient hypotheses and the meaning of maximal reduced part are not supplied; this assertion is retained as unresolved contextual information and is not used as a theorem or an additional hypothesis.
\subsection{Short English statement}
For the intended homogeneous spaces \(X=G/P\), with \(G\) simple adjoint and \(P\) a nonreduced parabolic subgroup scheme, characterize when the acting group equals the identity component of the automorphism group scheme, and describe that component when the supplied inclusion is strict. The base field, subgroup-scheme conventions, and faithfulness hypotheses require clarification before this becomes a fully specified problem.
\subsection{Sources}
\begin{itemize}
\item Supplied local UnsolvedMath registry record, permanent ID \(30006260\), source problem code \texttt{catalogue entry}, titled \emph{Acting Groups and Identity-Component Automorphisms}, in the Oberwolfach Reports Open Problems collection. The record gives proposed year 2025 and supplies no individual proposer.
\item Supplied report citation: \emph{Algebraic Groups} (2025), Oberwolfach Reports, DOI 10.4171/owr/2025/20, pp.\ 941--1010. \url{https://doi.org/10.4171/owr/2025/20}. The supplied background describes these reports as documenting mathematical workshops and their open-problem discussions.
\item Supplied literature source, labeled ``OWR 2025 report'': \url{https://ems.press/content/serial-article-files/51849}.
\item Supplied literature source, labeled ``Brion--Schroeer, inverse Galois problem'': \url{https://arxiv.org/abs/2205.08117}. No precise theorem or passage connecting this reference to the stated criterion is supplied.
\item Literature assessment dated \texttt{2026-08-22}: the supplied curation labels the problem open, says the extracted fragment omits required setup, and states that a broad realization theorem does not settle the exact isomorphism criterion. This is supplied dated curation, not an independently verified literature assessment.
\item The registry labels the statement \texttt{corrected\_verified} and reports that the official report was checked, the mathematical question isolated, and its notation repaired. These are supplied verification claims; no independent source inspection or mathematical verification is asserted here.
\end{itemize}
\end{document}
